% carousel/carousel.tex — LinkedIn-Karussell zum Whitepaper
% „Agentic Software Development — Enterprise Architecture with AI Agents"
% (Martin Janda, Fassung 2.1, August 2026)
%
% Separates Dokument; das Whitepaper (main.tex) bleibt unverändert.
% Eine Seite = eine Folie, quadratisch 1080 bp × 1080 bp (1080×1080-px-Äquivalent).
% Build:  latexmk -pdf -interaction=nonstopmode carousel.tex   (im Ordner carousel/)
%         bzw. `make karussell` im Projektwurzelverzeichnis.
%
% INHALTLICHE HERKUNFT: Jede Folie trägt einen Kommentar „QUELLE:", der auf
% Kapitel/Abschnitt der Whitepaper-Quelltexte (kapitel/*.md) verweist.
% Es sind ausschließlich Aussagen aus dem Whitepaper übernommen, nichts erfunden.
% Redigierbare Reintexte aller Folien: carousel/folientexte.txt
%
% ABGRENZUNG IMPLEMENTIERT / GEPLANT: Das Whitepaper trennt konsequent zwischen
% implementiertem Stand und Zielarchitektur. Die Folien halten diese Trennung
% ein — Folie 9 ist ausdrücklich als geplant gekennzeichnet.

\documentclass[12pt]{article}

% Gemeinsamer Stil aller Karussells (Format, Fonts, Farben, Folien-Bausteine):
% karussell.sty ist identisch zu dem der DeFi-Publikation.
\usepackage{karussell}
\renewcommand{\gesamt}{10}

% --- Folieneigene Grafiken -------------------------------------------------
% Die Abbildungen des Whitepapers sind für 16 cm Satzspiegel gezeichnet und in
% einer quadratischen Folie nicht mehr lesbar. Die beiden Grafiken hier sind
% vereinfachte, folientaugliche Ableitungen davon (Quelle je Folie im
% Kommentar) — gleiche Farbrollen wie die Buchabbildungen: schwarze Rahmen,
% pastellblaue Flächen, blaue Pfeile.
\usepackage{tikz}
\usetikzlibrary{arrows.meta,positioning}
\definecolor{kasten}{HTML}{DBE9F8}
\tikzset{
  kn/.style   = {draw=tinte, line width=1.1pt, fill=kasten, rounded corners=6pt,
                 align=center, inner sep=14pt, font=\fontsize{27}{33}\selectfont},
  knw/.style  = {kn, fill=white},
  pf/.style   = {-{Stealth[length=13pt,width=9pt]}, draw=blau, line width=2.2pt},
}

\begin{document}

% ============================================================================
% FOLIE 1 — Hook
% QUELLE: Kapitel 1, Absatz 1: Coding-Agenten nutzen Werkzeuge selbstständig,
% was ihre Ausführung „nichtdeterministisch und sicherheitsrelevant" macht;
% „Für den Enterprise-Einsatz genügt daher nicht die Auswahl eines
% leistungsfähigen Modells. Erforderlich ist eine Prozess- und
% Governance-Schicht, die Agentenarbeit spezifiziert, begrenzt, isoliert,
% prüft und nachweisbar macht."
% ============================================================================
\begin{folie}{Agentic Software Development}
  \gross{Der Engpass ist nicht das Modell. Es ist die fehlende
    Prozessschicht.}
  \vspace{48pt}
  \erl{Coding-Agenten rufen Werkzeuge selbstständig auf. Das macht sie
    produktiv — und ihre Ausführung nichtdeterministisch und
    sicherheitsrelevant. Befunde der überarbeiteten Fassung, Stand
    August 2026.}
\end{folie}

% ============================================================================
% FOLIE 2 — Kernbefund: Control Plane
% QUELLE: Management Summary, Kernaussage 1: „Der entscheidende Schritt ist
% nicht die maximale Zahl gleichzeitig arbeitender Agenten, sondern eine
% Steuerschicht, die nichtdeterministische Agentenarbeit begrenzt, isoliert,
% koordiniert, prüft und nachweisbar macht."
% Grafik: abb21 (Systemkontext der Control Plane), Abschnitt 19.1.
% ============================================================================
\begin{folie}{Der Kernbefund}
  \grossM{Nicht möglichst viele Agenten. Eine Steuerschicht dazwischen.}
  \vspace{58pt}
  {\centering
  \begin{tikzpicture}[node distance=30pt, scale=1, transform shape]
    \node[knw] (mensch) {Mensch\\[-2pt]{\fontsize{21}{25}\selectfont\color{tintesek}spezifiziert, gibt frei}};
    \node[kn, right=of mensch] (cp) {\textbf{Control Plane}\\[-2pt]{\fontsize{21}{25}\selectfont\color{tintesek}begrenzt, isoliert, prüft}};
    \node[knw, right=of cp] (agent) {Coding-Agent\\[-2pt]{\fontsize{21}{25}\selectfont\color{tintesek}schreibt}};
    \node[knw, below=42pt of cp] (repo) {Repository \textperiodcentered{} Audit-Kette};
    \draw[pf] (mensch) -- (cp);
    \draw[pf] (cp) -- (agent);
    \draw[pf] (cp) -- (repo);
  \end{tikzpicture}\par}
  \vspace{58pt}
  \erl{Die Control Plane sitzt zwischen Mensch, Coding-Agent und
    Zielrepository. Sie begrenzt, isoliert, koordiniert, prüft und macht
    nachweisbar.}
\end{folie}

% ============================================================================
% FOLIE 3 — Single Writer, Multiple Reviewers
% QUELLE: Prinzip AP-2 (Kapitel 2): „Pro Branch, Worktree oder Workspace
% schreibt zu einem Zeitpunkt genau ein Agent. Das verhindert
% Schreibkonkurrenz, unklare Zurechnung und nicht attestierbare
% Zwischenstände." Management Summary, Kernaussage 2.
% ============================================================================
\begin{folie}{Single Writer}
  \gross{Ein Schreiber je Arbeitskopie. Beliebig viele Prüfer.}
  \vspace{48pt}
  \erl{Zwei Agenten, die gleichzeitig dieselbe Arbeitskopie verändern,
    erzeugen Schreibkonkurrenz, unklare Zurechnung und Zwischenstände, für
    die niemand mehr geradestehen kann.}
\end{folie}

% ============================================================================
% FOLIE 4 — Independent Verification
% QUELLE: Prinzip AP-4 (Kapitel 2), Kurzformel „Wer schreibt, gibt nicht
% frei": „Erzeugung und Freigabe sind technisch getrennt: Read-only-Reviewer
% und Gates liefern strukturierte Befunde; ein Agent bewertet nie allein die
% eigene Arbeit."
% ============================================================================
\begin{folie}{Unabhängige Prüfung}
  \gross{Wer schreibt, gibt nicht frei.}
  \vspace{48pt}
  \erl{Erzeugung und Freigabe sind technisch getrennt. Read-only-Reviewer
    liefern strukturierte Befunde — und ein Agent bewertet nie allein die
    eigene Arbeit.}
\end{folie}

% ============================================================================
% FOLIE 5 — Fail Closed
% QUELLE: Prinzip AP-7 (Kapitel 2): „Ein ausgefallener Pflicht-Reviewer, eine
% nicht prüfbare Policy oder eine fehlgeschlagene Attestierung darf niemals
% als Erfolg gelten. Ein Gate, dessen Ausfall wie Erfolg aussieht, ist
% schlimmer als kein Gate."
% ============================================================================
\begin{folie}{Fail Closed}
  \grossM{Ein Gate, dessen Ausfall wie Erfolg aussieht, ist schlimmer als
    kein Gate.}
  \vspace{48pt}
  \erl{Ein ausgefallener Pflicht-Reviewer, eine nicht prüfbare Policy, eine
    fehlgeschlagene Attestierung: Nichts davon darf jemals als bestanden
    durchgehen.}
\end{folie}

% ============================================================================
% FOLIE 6 — Zustandshoheit
% QUELLE: Management Summary, Kernaussage 5: „Git trägt Artefakte, Branches
% und Checkpoints; Workflow-, Policy-, Freigabe- und Audit-Zustand liegen
% autoritativ in der Datenbank." Vgl. Kapitel 20, ADR-8.
% ============================================================================
\begin{folie}{Zustandshoheit}
  \gross{Git trägt die Artefakte. Den Prozesszustand trägt es nicht.}
  \vspace{48pt}
  \erl{Branches und Checkpoints gehören in Git. Workflow-, Policy-,
    Freigabe- und Audit-Zustand brauchen eine autoritative Quelle
    daneben — sonst ist der Prozess nicht rekonstruierbar.}
\end{folie}

% ============================================================================
% FOLIE 7 — Nachweisführung
% QUELLE: Abschnitt 19.1, Schritt 6 („Belegen"): „Jeder relevante Schritt
% landet in einer signierten, append-only Audit-Hashkette; ein Warum-Trace
% rekonstruiert für jeden Run, welche Policy, welches Modell und welche
% Freigabe gewirkt haben." Grafik: abb28 (Abschnitt 19.6).
% ============================================================================
\begin{folie}{Nachweis}
  \grossM{Für jeden Lauf rekonstruierbar: welche Policy, welches Modell,
    welche Freigabe.}
  \vspace{58pt}
  {\centering
  \begin{tikzpicture}[node distance=26pt]
    \node[knw] (pol) {Policy\\{\fontsize{21}{25}\selectfont\color{tintesek}versioniert}};
    \node[kn, right=of pol] (enf) {Durchsetzung\\{\fontsize{21}{25}\selectfont\color{tintesek}im Orchestrator}};
    \node[kn, right=of enf] (glied) {Kettenglied\\{\fontsize{21}{25}\selectfont\color{tintesek}signiert}};
    \node[knw, right=of glied] (exp) {Audit-Export\\{\fontsize{21}{25}\selectfont\color{tintesek}prüfbar}};
    \draw[pf] (pol) -- (enf);
    \draw[pf] (enf) -- (glied);
    \draw[pf] (glied) -- (exp);
  \end{tikzpicture}\par}
  \vspace{58pt}
  \erl{Jede Durchsetzung erzeugt ein signiertes Kettenglied in einer
    append-only Audit-Hashkette. Governance ist damit keine Ergänzung der
    Architektur, sondern Teil ihrer Struktur.}
\end{folie}

% ============================================================================
% FOLIE 8 — Referenzimplementierung
% QUELLE: Abschnitt 19, „Kennzahlen der Codebasis", erhoben auf dem
% Entwicklungsstand vom 26. Juli 2026: 362 Produktivklassen, ~33.000 Zeilen
% Produktivcode, 27 fachliche Slices, 26 Releases (0.1.0 bis 0.19.0,
% April–Juli 2026), Coverage-Gate Line >= 85 % / Branch >= 81 %
% (JaCoCo, buildbrechend).
% ============================================================================
\begin{folie}{Referenzimplementierung}
  \gross{Kein Konzeptpapier: 27 fachliche Slices, 26 Releases in vier
    Monaten.}
  \vspace{48pt}
  \erl{Rund 33.000 Zeilen Produktivcode, Testabdeckung als buildbrechendes
    Gate (Line ab 85 \%, Branch ab 81 \%). Die Architektur des Whitepapers
    war die Grundlage — das Kapitel beschreibt, was daraus geworden ist.
    Stand: 26.\,Juli 2026.}
\end{folie}

% ============================================================================
% FOLIE 9 — Roadmap (ausdrücklich als geplant gekennzeichnet)
% QUELLE: Management Summary, Absatz 4, und Abschnitt 19.10: Ein Feature wird
% in einen Task Graph zerlegt; unabhängige Tasks laufen als isolierte Child
% Runs; „Jeder Child Run besitzt einen eigenen Branch oder Worktree und bleibt
% dem Single-Writer-Prinzip unterworfen. Ein Merge Coordinator und ein
% abschließendes Integration Gate führen die Ergebnisse kontrolliert zusammen."
% Prinzip AP-3: Parallelität wird aus Abhängigkeiten abgeleitet, nicht aus
% festen Agentenrollen.
% ============================================================================
\begin{folie}{Nächster Schritt (geplant)}
  \grossM{Parallelität entsteht aus Abhängigkeiten — nicht aus
    Agentenrollen.}
  \vspace{48pt}
  \erl{Ein Feature wird in einen Task Graph zerlegt; unabhängige Tasks
    laufen als isolierte Child Runs, jeder mit eigenem Worktree und weiterhin
    mit genau einem Schreiber. Merge Coordinator und Integration Gate führen
    zusammen. Das Whitepaper trennt diesen Zielstand konsequent vom heute
    implementierten Stand.}
\end{folie}

% ============================================================================
% FOLIE 10 — Abschluss
% QUELLE: Titel und Untertitel der Arbeit (main.tex); Veröffentlichung auf
% janda.io. Kein QR-Code, keine URL-Parameter (Vorgabe wie bei Karussell 1).
% ============================================================================
\begin{folie}{Whitepaper}
  \grossM{Agentic Software Development}
  \vspace{40pt}
  {\fontsize{34}{46}\selectfont\color{tintesek}Enterprise Architecture with
    AI Agents — eine Control-Plane-Architektur für den kontrollierten Einsatz
    von Coding-Agenten\par}
  \vspace{64pt}
  \erl{\textcolor{tinte}{\bfseries Vollständiges Whitepaper auf janda.io}\\
    Überarbeitete und erweiterte Fassung 2.1, Stand: August 2026.}
\end{folie}

\end{document}
